\documentclass[11pt,letterpaper]{article}

% ── Packages ──────────────────────────────────────────────────────────────
\usepackage[margin=1in]{geometry}
\usepackage{titlesec}
\usepackage{enumitem}
\usepackage{hyperref}
\usepackage{booktabs}
\usepackage{longtable}
\usepackage{array}
\usepackage{xcolor}
\usepackage{fancyhdr}
\usepackage{parskip}
\usepackage{graphicx}

% ── Header / Footer ──────────────────────────────────────────────────────
\pagestyle{fancy}
\fancyhf{}
\rhead{\small dbGaP Data Access Request}
\lhead{\small Singlet AI --- Grand Rapids, MI}
\cfoot{\thepage}

\hypersetup{
    colorlinks=true,
    linkcolor=blue!60!black,
    urlcolor=blue!60!black,
}

\titleformat{\section}{\large\bfseries}{\thesection.}{0.5em}{}
\titleformat{\subsection}{\normalsize\bfseries}{\thesubsection.}{0.5em}{}

% ── Title ─────────────────────────────────────────────────────────────────
\title{%
    \vspace{-1.5cm}
    \textbf{dbGaP Data Access Request}\\[0.3em]
    \large Comprehensive Single-Cell Atlas for Non-Negative Matrix\\
    Factorization-Based Transcriptomic Reference Construction\\[0.3em]
    \normalsize Principal Investigator: Zach DeBruine, PhD
}
\author{
    Singlet AI\\
    Grand Rapids, Michigan\\
    \texttt{zach@singlet.ai}
}
\date{\today}

\begin{document}
\maketitle
\thispagestyle{fancy}

% ══════════════════════════════════════════════════════════════════════════
\section{Project Description}
% ══════════════════════════════════════════════════════════════════════════

\subsection{Background and Significance}

Single-cell and single-nucleus RNA sequencing (scRNA-seq / snRNA-seq) has
transformed our understanding of cellular heterogeneity in health and
disease. Publicly deposited datasets in NCBI GEO now exceed 180,000
samples, yet a substantial fraction of the most clinically valuable
single-cell data---particularly from studies involving human subjects
with protected health information---resides in controlled-access
repositories under dbGaP.

We have developed a scalable pipeline (\texttt{scgeo}) that discovers,
reprocesses, and harmonizes single-cell FASTQ data into
splicing-resolved count matrices using a uniform reference and
quantification workflow (STARsolo / Cell Ranger / Salmon alevin-fry).
The harmonized count matrices are then decomposed via
\textbf{Non-Negative Matrix Factorization (NMF)} into a genes-by-factors
matrix $\mathbf{W}$ and a factors-by-cells matrix $\mathbf{H}$, which
together define a universal transcriptomic reference that captures
conserved gene programs across tissues, diseases, and developmental
contexts.

Incorporating the controlled-access data in this reference is critical
because:
\begin{itemize}[nosep]
    \item Disease-specific cohorts (cancer, neurodegeneration,
          autoimmune) are disproportionately represented in dbGaP
          relative to GEO.
    \item Paired whole-genome sequencing (WGS) and single-cell
          transcriptome data---available in studies such as HTAN, GTEx,
          and BICCN---enable linking genomic variation to
          transcriptomic cell states, a key requirement for integration
          with emerging foundation models (e.g., Google AlphaGenome).
    \item Atlas-scale programs (HuBMAP, KPMP, SenNet) provide
          comprehensive healthy-tissue references that anchor
          disease-state comparisons.
\end{itemize}

\subsection{Research Use Statement}

The requested individual-level data will be used exclusively for the
following non-commercial, academic research purpose:

\begin{quote}
\textit{To construct a comprehensive, splicing-resolved single-cell
transcriptomic reference by reprocessing raw FASTQ files from
controlled-access single-cell and single-nucleus sequencing studies
into harmonized count matrices, and decomposing those matrices via
Non-Negative Matrix Factorization (NMF) to identify universal gene
expression programs. The resulting factor model
($\mathbf{W} \in \mathbb{R}^{G \times K}$,
 $\mathbf{H} \in \mathbb{R}^{K \times N}$)
captures population-level gene programs without retaining
individual-level identifiable information. Only the $\mathbf{W}$ matrix
(gene loadings) and aggregated per-factor summary statistics will be
distributed; individual-level $\mathbf{H}$ coefficients, raw reads,
and donor metadata will remain within the approved secure compute
environment.}
\end{quote}

\subsection{Specific Aims}

\begin{enumerate}[nosep]
    \item \textbf{Aim 1: Harmonized Reprocessing.}
          Download raw FASTQ files for all requested studies to the
          approved university HPC cluster. Align reads to a uniform
          reference genome (GRCh38 / GRCm39) using STARsolo with
          splice-aware quantification, producing per-sample count
          matrices with spliced / unspliced layers.

    \item \textbf{Aim 2: NMF Decomposition.}
          Integrate the reprocessed count matrices with our existing
          public-data catalog ($\sim$180K samples from GEO + 894
          experiments from ENCODE) and perform rank-optimized NMF to
          derive a universal gene-program dictionary $\mathbf{W}$.

    \item \textbf{Aim 3: Genome--Transcriptome Linkage.}
          For studies with paired WGS/WES and single-cell
          transcriptome data (e.g., HTAN, GTEx, BICCN), project
          individual cells onto the NMF factor space and associate
          factor usage patterns with germline and somatic variants to
          enable integration with genomic foundation models.

    \item \textbf{Aim 4: Public Reference Release.}
          Release only the non-identifiable $\mathbf{W}$ matrix and
          aggregated factor-level summary statistics (k-anonymized,
          $n \geq 10$ donors per group) as a public transcriptomic
          reference. No individual-level data, $\mathbf{H}$ matrix
          values, or donor metadata will be redistributed.
\end{enumerate}

% ══════════════════════════════════════════════════════════════════════════
\section{Requested Datasets}
% ══════════════════════════════════════════════════════════════════════════

We request access to the following 121 dbGaP studies, grouped by
consortium and disease focus. All studies contain single-cell or
single-nucleus sequencing data with raw FASTQ files deposited in SRA
under controlled access.

\subsection{Major Consortium Studies}

\begin{longtable}{@{} l l p{7.5cm} @{}}
\toprule
\textbf{phs Accession} & \textbf{Consortium} & \textbf{Description} \\
\midrule
\endhead
phs002371 & HTAN & Human Tumor Atlas Network --- pan-cancer SC + spatial + WGS \\
phs000424 & GTEx & Genotype-Tissue Expression --- snRNA-seq + WGS across 8 tissues \\
phs002673 & BICCN & Brain Initiative Cell Census Network --- 50M+ brain cells \\
phs003431 & BICAN & Brain Initiative Cell Atlas Network --- next-gen brain atlas \\
phs003604 & BICAN-2 & BICAN Phase 2 \\
phs001810 & HuBMAP & Human BioMolecular Atlas Program --- SC across organs \\
phs002507 & HuBMAP-2 & HuBMAP Phase 2 \\
phs000178 & GDC/TCGA & The Cancer Genome Atlas --- WGS/WES + SC \\
phs000218 & TARGET & Therapeutically Applicable Research --- pediatric cancers \\
phs001211 & TOPMed & Trans-Omics for Precision Medicine --- WGS + phenotypes \\
phs002362 & HTAN-HTAPP & Human Tumor Atlas Pilot --- SC + WGS \\
phs002624 & SenNet & Cellular Senescence Network --- SC atlas of aging \\
phs002256 & KPMP & Kidney Precision Medicine Project --- snRNA + spatial \\
phs003169 & HCA & Human Cell Atlas --- large-scale SC reference \\
phs002716 & AMP-PD & AMP Parkinson's Disease --- WGS + snRNA \\
phs001515 & LungMAP & NHLBI Lung Molecular Atlas Program \\
phs002062 & 4DN & 4D Nucleome --- chromatin architecture \\
phs003366 & IGVF & Impact of Genomic Variation on Function \\
phs003123 & SPARC & Stimulating Peripheral Activity \\
phs002159 & HPAP & Human Pancreas Analysis Program \\
phs002279 & ABI & Allen Brain Institute Cell Types \\
phs001246 & SEA-AD & Seattle Alzheimer's Disease Brain Cell Atlas \\
\bottomrule
\end{longtable}

\subsection{Individual Investigator Studies}

The remaining 99 phs accessions represent individual investigator-led
studies spanning cancer (36 studies), neurological disease (28),
immune/autoimmune (15), developmental biology (10), and other categories
(10). A complete list with phs accessions and titles is provided in
\textbf{Supplementary Table S1} appended to this application.

\subsection{Data Types Requested}

For each study we request:
\begin{itemize}[nosep]
    \item \textbf{Raw sequencing data (FASTQ/BAM):} Required for
          uniform reprocessing with our harmonized pipeline.
    \item \textbf{Sample and donor metadata:} Cell type annotations,
          tissue of origin, disease status, age, sex, treatment---needed
          for meaningful NMF factor annotation. All metadata remains
          within the secure compute environment.
    \item \textbf{Genotype/WGS data (where available):} For the 153
          studies with paired WGS/WES data, used exclusively for
          genome--transcriptome linkage analysis (Aim~3).
\end{itemize}

% ══════════════════════════════════════════════════════════════════════════
\section{Data Security and Institutional Oversight}
% ══════════════════════════════════════════════════════════════════════════

\subsection{Institutional Affiliation}

\begin{tabular}{@{} l l @{}}
\textbf{Institution:} & \textit{[University Name]} \\
\textbf{Signing Official:} & \textit{[Name, Title]} \\
\textbf{IT Security Officer:} & \textit{[Name, Title]} \\
\textbf{IRB of Record:} & \textit{[IRB Protocol Number]} \\
\end{tabular}

\bigskip
\noindent The Principal Investigator has completed the following:
\begin{itemize}[nosep]
    \item NIH Genomic Data Sharing (GDS) Responsible Conduct of Research training
    \item CITI Human Subjects Research training
    \item Institutional Security Awareness training
\end{itemize}

\subsection{IT Security Plan}

Data will be stored and processed on the university's
\textbf{High-Performance Computing (HPC)} cluster, which meets or
exceeds the NIH Security Best Practices for Controlled-Access Data:

\begin{longtable}{@{} p{4cm} p{9cm} @{}}
\toprule
\textbf{Requirement} & \textbf{Implementation} \\
\midrule
\endhead
Physical security & University data center with badge access, 24/7 surveillance \\
Network security & Campus firewall, intrusion detection (IDS/IPS), VPN-only access \\
Authentication & University SSO with MFA; SSH key-based access to HPC \\
Authorization & POSIX ACLs restricting data directories to PI + approved personnel \\
Encryption at rest & AES-256 on GPFS/Lustre filesystem or encrypted scratch \\
Encryption in transit & TLS 1.3 for all data transfers; SCP/SFTP for file movement \\
Logging \& auditing & All access logged via \texttt{auditd}; quarterly review \\
Data destruction & Secure deletion (NIST 800-88) upon project completion or expiry \\
Incident response & University CISO incident response plan; 24-hour breach notification \\
\bottomrule
\end{longtable}

\subsection{Data Use Limitations}

\begin{itemize}[nosep]
    \item \textbf{No redistribution} of individual-level data (FASTQ,
          BAM, count matrices, metadata).
    \item \textbf{No re-identification attempts} of study participants.
    \item \textbf{No cloud processing} --- all data remains on
          university on-premises HPC infrastructure.
    \item \textbf{Outputs released publicly} are limited to the
          population-level $\mathbf{W}$ matrix and k-anonymized
          ($k \geq 10$) summary statistics.
    \item Data access will be terminated and data securely destroyed
          upon project completion, change of PI affiliation, or
          revocation of dbGaP approval.
\end{itemize}

% ══════════════════════════════════════════════════════════════════════════
\section{Compute Architecture and Data Flow}
% ══════════════════════════════════════════════════════════════════════════

\subsection{Processing Pipeline}

\begin{enumerate}[nosep]
    \item \textbf{Download:} \texttt{prefetch} + \texttt{fasterq-dump}
          from SRA to university HPC scratch storage via authenticated
          dbGaP download (NGC token).
    \item \textbf{Alignment:} STARsolo (10x Chromium, Drop-seq,
          inDrop) or Salmon alevin-fry (Smart-seq2/3) on compute nodes.
          SLURM job scheduler, 64--128 cores per job.
    \item \textbf{QC \& Filtering:} Per-cell QC metrics (UMI counts,
          gene counts, mitochondrial fraction); doublet detection
          (Scrublet).
    \item \textbf{Count matrix output:} Sparse matrices (genes $\times$
          cells) with spliced/unspliced layers stored in
          \texttt{.h5ad} format on encrypted storage.
    \item \textbf{NMF training:} GPU-accelerated rank-optimized NMF on
          the merged count matrix. The $\mathbf{W}$ matrix is exported;
          $\mathbf{H}$ remains on secure storage.
    \item \textbf{Data destruction:} Raw FASTQ and intermediate files
          deleted upon successful QC of count matrices. Count matrices
          retained only for the duration of the approved project period.
\end{enumerate}

\subsection{Data Flow Diagram}

\begin{verbatim}
    SRA (dbGaP)                  University HPC (Secure)
    ┌──────────┐    NGC token    ┌────────────────────────────┐
    │  FASTQ   │ ──────────────> │  /scratch/dbgap/{phs}/     │
    │  (SRA)   │   prefetch +    │     ├── raw/  (FASTQ)      │
    └──────────┘  fasterq-dump   │     ├── aligned/ (BAM)     │
                                 │     ├── counts/ (h5ad)     │
                                 │     └── metadata/          │
                                 │                            │
                                 │  STARsolo / alevin-fry     │
                                 │         ↓                  │
                                 │  NMF (GPU nodes)           │
                                 │    W matrix  ──────────>  PUBLIC
                                 │    H matrix  (stays)       │
                                 │    metadata  (stays)       │
                                 └────────────────────────────┘
\end{verbatim}

% ══════════════════════════════════════════════════════════════════════════
\section{PHI and Privacy Considerations}
% ══════════════════════════════════════════════════════════════════════════

\subsection{Output Privacy Analysis}

The sole public output of this project is the
$\mathbf{W} \in \mathbb{R}^{G \times K}$ matrix, where $G \approx$
33,000 genes and $K \approx$ 100--500 factors. This matrix represents
\textbf{population-level gene program definitions} analogous to a
dictionary of gene co-expression patterns. It contains:
\begin{itemize}[nosep]
    \item No individual-level data
    \item No donor identifiers or demographic information
    \item No cell-level expression values
    \item No sequence-level information (no reads, no variants)
\end{itemize}

The $\mathbf{H}$ matrix (factors $\times$ cells) and all donor metadata
remain within the secure compute environment and are never distributed.

\textbf{Summary statistics} derived from $\mathbf{H}$ are released only
after k-anonymization: for each NMF factor, we report mean and variance
of factor usage across coarse-grained groups (tissue type, broad disease
category) with a minimum group size of $k = 10$ donors. This satisfies
standard k-anonymity requirements and prevents re-identification.

\subsection{Re-identification Risk Assessment}

The $\mathbf{W}$ matrix is a \textbf{generative model parameter} rather
than a transformation of individual records. Formally, $\mathbf{W}$ is
the solution to:
\[
    \min_{\mathbf{W} \geq 0,\, \mathbf{H} \geq 0}
    \| \mathbf{X} - \mathbf{W}\mathbf{H} \|_F^2
    + \lambda \|\mathbf{H}\|_1
\]
where $\mathbf{X}$ is the full count matrix. Since $\mathbf{W}$ is
determined by the global structure of $\mathbf{X}$ rather than any
individual column, removing or adding a single donor changes
$\mathbf{W}$ by $O(1/N)$ where $N$ is on the order of millions of
cells. Individual contribution is negligible and non-recoverable.

% ══════════════════════════════════════════════════════════════════════════
\section{Personnel}
% ══════════════════════════════════════════════════════════════════════════

\begin{longtable}{@{} l l p{5.5cm} l @{}}
\toprule
\textbf{Name} & \textbf{Role} & \textbf{Responsibilities} & \textbf{Training} \\
\midrule
\endhead
Zach DeBruine, PhD & PI & Study design, NMF methods, oversight & GDS, CITI \\
\textit{[TBD]} & Co-I & Pipeline development, QC & GDS, CITI \\
\textit{[TBD]} & Data Analyst & Reprocessing, metadata & GDS, CITI \\
\bottomrule
\end{longtable}

All personnel with data access will complete required training prior to
data download. Access will be revoked immediately upon departure from the
approved institution.

% ══════════════════════════════════════════════════════════════════════════
\section{Timeline}
% ══════════════════════════════════════════════════════════════════════════

\begin{tabular}{@{} l l @{}}
\toprule
\textbf{Milestone} & \textbf{Target} \\
\midrule
DAR approval & $T_0$ \\
Infrastructure validation \& test downloads & $T_0$ + 2 weeks \\
Phase 1: Consortium studies (HTAN, GTEx, BICCN, HuBMAP) & $T_0$ + 3 months \\
Phase 2: Disease-specific studies (cancer, neuro) & $T_0$ + 6 months \\
Phase 3: Remaining studies + paired WGS analysis & $T_0$ + 9 months \\
NMF model training and validation & $T_0$ + 12 months \\
Data destruction (raw files) & $T_0$ + 14 months \\
Annual renewal / closeout & $T_0$ + 12 months \\
\bottomrule
\end{tabular}

% ══════════════════════════════════════════════════════════════════════════
\section*{Supplementary Table S1: Complete phs Accession List}
% ══════════════════════════════════════════════════════════════════════════
\addcontentsline{toc}{section}{Supplementary Table S1}

{\small
\begin{longtable}{@{} l p{10cm} @{}}
\toprule
\textbf{phs Accession} & \textbf{Study Title} \\
\midrule
\endhead
phs000007 & NHLBI Framingham SABRe CVD \\
phs000159 & Therapeutic targets in acute myeloid leukemia \\
phs000178 & The Cancer Genome Atlas (GDC/TCGA) \\
phs000187 & High Density SNP Association Analysis of Melanoma \\
phs000218 & Therapeutically Applicable Research (TARGET) \\
phs000363 & NHLBI Framingham SABRe CVD \\
phs000424 & Genotype-Tissue Expression (GTEx) \\
phs000549 & Clonal Evolution of Pre-Leukemic HSC \\
phs000640 & Massively Multiplex Single-Cell Hi-C \\
phs000748 & Therapeutic targets --- multi-omic characterization \\
phs000815 & Gene Expression and Regulatory Networks in Leukocytes \\
phs000833 & Single Cell Analysis Program -- Transcriptome (SCAP-T) \\
phs000834 & SCAP-T (UCSD) \\
phs000835 & SCAP-T (U.\ Penn) \\
phs000836 & SCAP-T (USC) \\
phs000924 & iPSCORE (iPSC Collection for Omic Research) \\
phs000933 & Acral melanoma molecular profiling \\
phs000958 & CSER-MedSeq \\
phs000989 & Human Developmental Neurogenesis \\
phs000993 & GTEx Pilot \\
phs001016 & Fixed Single Cell Study (Radial Glia) \\
phs001024 & NHGRI GREGoR \\
phs001027 & SC genomics --- Homo sapiens \\
phs001029 & Full-Length scRNA-seq (HeLa S3) \\
phs001048 & FUSION Tissue Biopsy Study \\
phs001068 & FUSION Muscle Expression \\
phs001110 & Kids First --- Congenital Heart Disease \\
phs001158 & Spatial transcriptomics of DRG (pain) \\
phs001181 & SC DNA and RNA Sequencing of CLL \\
phs001187 & Epigenomics of CD8 T cell Aging \\
phs001188 & FUSION Islet Expression \\
phs001200 & West Nile Virus Neutralizing Antibodies \\
phs001205 & hESC-derived brain cell lineages \\
phs001211 & TOPMed \\
phs001231 & NeuroLINCS \\
phs001246 & Seattle Alzheimer's Disease (SEA-AD) \\
phs001255 & Molecular Evolution of Cancer \\
phs001268 & Combinatorial Indexing SC Genomes \\
phs001269 & Massively Multiplex SC Hi-C \\
phs001276 & hESC-derived interneuron development \\
phs001294 & SC genomics study \\
phs001300 & Single Nuclei Multiome -- NABEC Cortex \\
phs001323 & Immune biomarkers of checkpoint response \\
phs001362 & Linked Read Sequencing -- Gastric Cancer \\
phs001372 & Integrated SC Analysis in CLL \\
phs001378 & Lymphoma scRNA-Seq \\
phs001438 & Chromatin Accessibility in Developing Brain \\
phs001457 & AMP-RA (Rheumatoid Arthritis) \\
phs001515 & NHLBI LungMAP \\
phs001519 & Neoantigen Vaccine in GBM \\
phs001525 & X-Linked Dystonia Parkinsonism \\
phs001529 & scRNA-Seq of RA Synovial Tissue \\
phs001543 & VAST --- Vascular Smooth Muscle \\
phs001587 & Cellular heterogeneity across lineage \\
phs001615 & Epigenetic Landscape of RA FLS \\
phs001627 & SC DNA Sequencing of AML \\
phs001671 & U6 snRNA / LINE-1 Retrotransposition \\
phs001680 & T Cell States and Checkpoint Blockade in Melanoma \\
phs001703 & Genetic Polymorphisms and Immune Gene Expression \\
phs001704 & Enhancer Mapping in CLL \\
phs001750 & IPF Cell Atlas \\
phs001763 & All of Us Research Program \\
phs001773 & Linked-Read Sequencing --- Complex SVs \\
phs001790 & NIMH Human MTG Cell Types \\
phs001810 & HuBMAP \\
phs001836 & SC Transcriptomics of Neocortical Development \\
phs001840 & SC Analysis of Pancreatic Ductal Adenocarcinoma \\
phs001854 & Pilocytic Astrocytoma scRNA-Seq \\
phs001861 & Uveal Melanoma SC Evolution \\
phs001886 & SC Analysis of Human Placenta \\
phs001907 & BRAIN Initiative SCAP-T Phase 2 \\
phs001934 & Sexual Dimorphism in Immune Aging \\
phs001961 & SN Multiomic Profiling of Host Response \\
phs002003 & Fetal Gene Expression Atlas \\
phs002007 & Innate T Cell Transcriptomics \\
phs002008 & Orofacial Morphology Enhancers \\
phs002015 & SCAP-T Phase 2 (Broad) \\
phs002048 & SC Analysis of Systemic Lupus \\
phs002062 & 4D Nucleome \\
phs002071 & Donor Pancreata Transcriptomic Signatures \\
phs002159 & Human Pancreas Analysis Program (HPAP) \\
phs002204 & Cardiac Cell Type Gene Regulatory Programs \\
phs002256 & KPMP \\
phs002258 & HEAL Initiative \\
phs002279 & Allen Brain Institute Cell Types \\
phs002280 & Multimodal Brain Analysis \\
phs002316 & Trimodal SC Measurement (DOGMA-seq) \\
phs002361 & Baseline Immune and Checkpoint Response \\
phs002362 & HTAN-HTAPP \\
phs002371 & HTAN \\
phs002386 & H3Africa \\
phs002407 & MMRd/MMRp Colorectal Cancer Atlas \\
phs002443 & In-situ to Invasive Breast Cancer \\
phs002507 & HuBMAP Phase 2 \\
phs002624 & SenNet \\
phs002655 & SARS-CoV-2 Immune Perturbations in Infants \\
phs002673 & BICCN \\
phs002716 & AMP-PD \\
phs002790 & NCI/IOTN Immunotherapy \\
phs002808 & NCI Moonshot --- Tumor Atlas (TNET) \\
phs002915 & SC Multi-omic Velocity --- Gene Regulation \\
phs002922 & Baseline Immune and T Cell Kinetics \\
phs002926 & Infant Serological Responses \\
phs003002 & CD137 Agonism + Anti-PD1 \\
phs003015 & Donor/Recipient Immune Phenotypes Post-Transplant \\
phs003123 & SPARC \\
phs003169 & Human Cell Atlas \\
phs003230 & Neuroendocrine Prostate Cancer Evolution \\
phs003260 & SN Transcriptomes --- Opioid Overdose \\
phs003366 & IGVF \\
phs003417 & Multimodal Gene Expression at Cell Level \\
phs003418 & SC Characterization of Microscopic Colitis \\
phs003427 & SC Profiling of Premature Neonate Airways \\
phs003431 & BICAN \\
phs003604 & BICAN Phase 2 \\
\bottomrule
\end{longtable}
}

\end{document}
